\documentclass[10pt]{article}
\usepackage[margin=0.7in]{geometry}
\usepackage{booktabs}
\usepackage{xcolor}
\usepackage{enumitem}
\usepackage{microtype}
\usepackage{titlesec}

\definecolor{accent}{HTML}{1F3C5A}
\titleformat{\section}{\normalsize\bfseries\color{accent}}{}{0em}{}
\titlespacing*{\section}{0pt}{0.45em}{0.25em}

\pagestyle{empty}
\setlength{\parindent}{0pt}
\setlength{\parskip}{0.2em}
\setlist[itemize]{leftmargin=1.1em,itemsep=0.05em,topsep=0.1em}

\begin{document}
\small

\begin{center}
{\Large\bfseries Root-Tweet Forward Prediction of Political Reply Dynamics}\\[0.25em]
{\normalsize Technical Whitepaper}\\[0.2em]
Luke Tervit \quad|\quad University of Edinburgh \quad|\quad March 2026
\end{center}

\section{Objective}
This project asks whether downstream political discourse can be predicted from the \textbf{root tweet alone}. The modelling goal is not exact next-token recreation of future replies, but distribution-level forecasting of how a conversation unfolds across political alignment, emotion, sentiment, and aggression.

\section{System Summary}
The pipeline combines an agent-based model (Mesa) with LLM generation (Dolphin-Llama3 8B via Ollama):
\begin{itemize}
  \item Agent personas are constructed from historical Twitter data using five classifiers: political leaning, emotion, sentiment, hate, and offensive language.
  \item Each simulation starts with only the root tweet in thread metadata (0-shot prediction mode; no real target replies exposed).
  \item Agents interact over 10 synchronous rounds; reply probability is behaviourally weighted (base + aggression boost), and each generated post conditions on recent thread context.
  \item Validation reclassifies real and simulated tweets with the same five-model stack and computes JSD-based similarity plus aggression-mean gap.
\end{itemize}

\section{Specific-User Forward Holdout Experiment}
\textbf{Target account:} Sen. Marsha Blackburn (user\_id 278145569).\\
\textbf{Protocol:} 8 chronological root tweets found; first 3 used to build responder personas; \textbf{last 5 held out} for forward prediction.\\
\textbf{Prediction input per run:} root tweet text only (\texttt{source: forward\_holdout\_zero\_shot}; \texttt{use\_few\_shot: false}).\\
\textbf{Population:} 200 sampled agents from 124 matched historical responders.

\begin{center}
\begin{tabular}{l c}
\toprule
Holdout Root Tweet ID & Overall Accuracy (\%) \\
\midrule
1800961685949198714 & 97.93 \\
1800996791078236518 & 94.83 \\
1801248635654390240 & 94.58 \\
1801290362700394984 & 90.57 \\
1801402597707767909 & 94.79 \\
\midrule
\textbf{Mean $\pm$ SD} & \textbf{94.54 $\pm$ 2.34} \\
\bottomrule
\end{tabular}
\end{center}

\textbf{Component means across the 5 predicted tweets:}
Political similarity 96.45\%, Emotion similarity 91.78\%, Sentiment similarity 97.39\%, Aggression similarity 90.17\%.

\section{Interpretation and Significance}
The core claim is supported: using only the root tweet and historically learned responder behaviour, the system forecasted the next five held-out conversations with \textbf{94.54\% mean weighted accuracy}. This extends broader project evidence (92.3\% mean on 100 held-out threads) by showing user-specific forward prediction works in strict chronological holdout mode.

\section{Boundaries and Next Steps}
\begin{itemize}
  \item The 94.54\% metric is distributional fidelity, not exact tweet-by-tweet text matching.
  \item Reply volume was not directly optimised in this score; future work should add calibrated volume/velocity objectives.
  \item External validity remains bounded to USC X 24 political discourse and tested model settings.
\end{itemize}

\vspace{0.25em}
\textbf{Bottom line:} past political tweets from a specific account can be transformed into a predictive behavioural population that anticipates likely reaction patterns to the account's next root tweets before those conversations are observed.

\end{document}
